\documentclass[11pt,a4paper]{article}

% --- Encoding ---
\usepackage[utf8]{inputenc}
\usepackage[T1]{fontenc}
\usepackage[spanish,es-noshorthands]{babel}
\usepackage{lmodern}

% --- Layout ---
\usepackage[margin=2.5cm]{geometry}
\usepackage{parskip}
\setlength{\parindent}{0pt}
\setlength{\parskip}{0.5em plus 0.2em minus 0.1em}

% --- Tables and graphics ---
\usepackage{booktabs}
\usepackage{array}
\usepackage{enumitem}
\usepackage{xcolor}
\usepackage{graphicx}
\usepackage{tikz}
\usetikzlibrary{positioning,arrows.meta}
\usepackage{caption}
\usepackage{subcaption}
\usepackage{amsmath,amssymb}

% --- Code listings ---
\usepackage{listings}
\definecolor{codegray}{rgb}{0.96,0.96,0.96}
\definecolor{codeborder}{rgb}{0.85,0.85,0.85}
\definecolor{codecomment}{rgb}{0.30,0.55,0.30}
\definecolor{codekeyword}{rgb}{0.10,0.30,0.70}
\definecolor{codestring}{rgb}{0.65,0.15,0.10}

\lstdefinestyle{shellstyle}{
    language=bash,
    backgroundcolor=\color{codegray},
    basicstyle=\ttfamily\small,
    breaklines=true,
    frame=single,
    rulecolor=\color{codeborder},
    showstringspaces=false,
    keywordstyle=\color{codekeyword}\bfseries,
    commentstyle=\color{codecomment}\itshape,
    stringstyle=\color{codestring},
    aboveskip=1em,
    belowskip=1em,
    columns=fullflexible,
}

\lstdefinestyle{cstyle}{
    language=C,
    backgroundcolor=\color{codegray},
    basicstyle=\ttfamily\footnotesize,
    breaklines=true,
    frame=single,
    rulecolor=\color{codeborder},
    showstringspaces=false,
    keywordstyle=\color{codekeyword}\bfseries,
    commentstyle=\color{codecomment}\itshape,
    stringstyle=\color{codestring},
    numbers=left,
    numberstyle=\tiny\color{gray},
    numbersep=8pt,
    aboveskip=1em,
    belowskip=1em,
}

\lstdefinestyle{javastyle}{
    language=Java,
    backgroundcolor=\color{codegray},
    basicstyle=\ttfamily\footnotesize,
    breaklines=true,
    frame=single,
    rulecolor=\color{codeborder},
    showstringspaces=false,
    keywordstyle=\color{codekeyword}\bfseries,
    commentstyle=\color{codecomment}\itshape,
    stringstyle=\color{codestring},
    numbers=left,
    numberstyle=\tiny\color{gray},
    numbersep=8pt,
    aboveskip=1em,
    belowskip=1em,
}

\lstdefinestyle{textstyle}{
    backgroundcolor=\color{codegray},
    basicstyle=\ttfamily\footnotesize,
    breaklines=true,
    frame=single,
    rulecolor=\color{codeborder},
    showstringspaces=false,
    aboveskip=1em,
    belowskip=1em,
    columns=fullflexible,
}

\lstset{style=shellstyle}

% --- Hyperlinks ---
\usepackage[colorlinks=true,
            linkcolor=blue!60!black,
            urlcolor=blue!50!black,
            citecolor=blue!60!black]{hyperref}

% --- Title block ---
\title{
    \vspace{-2.5em}
    \textbf{Práctico 3} \\[0.3em]
    \large Procesos, threads y concurrencia
}
\author{
    Tomás Rodeghiero \\
    \small Sistemas Operativos --- Universidad Nacional de Río Cuarto
}
\date{2026}

\begin{document}

\maketitle

\begin{abstract}
\noindent
Este documento presenta la resolución integral del Práctico 3 de la
materia Sistemas Operativos (UNRC, 2026), centrado en la administración
de \emph{procesos} y \emph{threads}, los mecanismos de \emph{planificación}
de CPU y los problemas de \emph{concurrencia} y \emph{sincronización}.
Cada ejercicio se justifica con la teoría de los capítulos
\emph{Procesos y threads} y \emph{Concurrencia e IPC} de las
\emph{Notas del curso} (Marcelo Arroyo) y se acompaña con comandos,
salidas, fragmentos de código y experimentos verificados sobre
GNU/Linux y macOS. La discusión recorre los estados de un proceso, el
control con señales, el árbol de procesos, la planificación con prioridades,
las condiciones de carrera, la sincronización con \texttt{flock}, mutexes
POSIX, monitores Java, las operaciones del kernel ante interrupciones de
timer y de I/O, la prueba semi-formal del algoritmo de Peterson y la
implementación del problema productor-consumidor con semáforos POSIX.
\end{abstract}

\tableofcontents
\bigskip

% =====================================================================
\section{Introducción}
% =====================================================================

Un \emph{proceso} es una instancia de un programa en ejecución; un
\emph{thread} (o \emph{tarea}) es la unidad mínima planificable por el
kernel. Un proceso tiene siempre al menos un thread, su \emph{main
thread}, y puede crear threads adicionales que comparten el espacio de
memoria del proceso pero tienen pilas independientes. Los recursos
asignados a un proceso ---memoria, descriptores de archivos, manejadores
de señales, áreas de memoria compartida--- son administrados
explícitamente por el kernel mediante una \emph{tabla de procesos}.

La teoría del curso modela el ciclo de vida de un proceso con cinco
estados: \texttt{NEW}, \texttt{RUNNING}, \texttt{RUNNABLE} (o
\texttt{READY}), \texttt{WAITING} (o \texttt{SLEEPING}), \texttt{ZOMBIE}
y \texttt{TERMINATED}. Las transiciones se disparan por
\emph{interrupciones}, \emph{traps} y \emph{llamadas al sistema}: el
timer fuerza la transición \texttt{RUNNING}~$\rightarrow$~\texttt{RUNNABLE}
al consumir el quantum; un \texttt{read} sobre disco produce
\texttt{RUNNING}~$\rightarrow$~\texttt{SLEEPING}; un \texttt{wakeup}
produce \texttt{SLEEPING}~$\rightarrow$~\texttt{RUNNABLE}; \texttt{exit}
lleva al proceso a \texttt{ZOMBIE} hasta que el padre haga \texttt{wait}.

Cuando varias tareas comparten un recurso aparece la \emph{concurrencia}.
La porción de código que accede al recurso compartido se llama
\emph{región crítica} y la dependencia del resultado en el orden de
ejecución se llama \emph{condición de carrera} (\emph{race condition}).
Para evitar trazas no deseadas el SO ofrece una jerarquía de
\emph{mecanismos de sincronización}: spinlocks, sleep locks, semáforos,
variables de condición, monitores y mensajes. Cada uno apunta a un
escenario distinto: regiones cortas vs. largas, monoprocesador vs.
multiprocesador, memoria compartida vs. procesos aislados.

El práctico ejercita progresivamente todos estos conceptos: parte del
\emph{observable} (estados de los procesos visibles con \texttt{ps},
control con señales, árbol con \texttt{pstree}), pasa por la
\emph{planificación con prioridades} (\texttt{nice}/\texttt{renice}),
introduce las \emph{race conditions} (\texttt{counter.c} con dos
procesos, \texttt{pthreads-example.c} con dos hilos) y sus correcciones
(\texttt{flock}, \texttt{pthread\_mutex\_t}, \texttt{synchronized} de
Java), describe el comportamiento del kernel en dos escenarios típicos
(timer y \texttt{read} con I/O), demuestra semi-formalmente el algoritmo
de Peterson y termina con el clásico problema del \emph{productor-consumidor}
implementado con semáforos POSIX, en sus dos variantes: con threads
(memoria compartida) y con procesos (named semaphores y shared memory).

% =====================================================================
\section{Ejercicio 1 --- Estado de los procesos con \texttt{ps}}
% =====================================================================

Se pide ejecutar \texttt{ps aux}, identificar los estados de los
procesos, ordenar por uso de CPU y de memoria, y listar los procesos
del usuario actual.

\subsection*{1.a Lectura de la salida y estados}

\begin{lstlisting}
ps aux | head -n 5
\end{lstlisting}

Cada fila representa una tarea. Las columnas claves son:
\texttt{USER} y \texttt{PID} (identificación), \texttt{\%CPU}/\texttt{\%MEM}
(consumo instantáneo), \texttt{VSZ}/\texttt{RSS} (memoria virtual y
residente), \texttt{TTY} (terminal controladora), \texttt{STAT} (estado
codificado), \texttt{TIME} (CPU acumulado) y \texttt{COMMAND} (línea de
comando).

La columna \texttt{STAT} expone los estados teóricos del modelo de
tareas. La correspondencia entre los códigos \texttt{ps} y los estados
del PCB es:

\begin{center}
\small
\begin{tabular}{@{}lllp{6.5cm}@{}}
\toprule
Estado teórico & Linux & macOS & Significado en el kernel \\
\midrule
\texttt{RUNNING}/\texttt{RUNNABLE} & \texttt{R} & \texttt{R} & En la CPU o en la cola de \emph{ready} \\
\texttt{SLEEPING} interrumpible & \texttt{S} & \texttt{S}/\texttt{I} & En \emph{wait queue}; despertable por señal \\
\texttt{SLEEPING} no interrumpible & \texttt{D} & \texttt{U} & I/O síncrono; no despertable por señal \\
Detenido por señal & \texttt{T} & \texttt{T} & \texttt{SIGSTOP}/\texttt{SIGTSTP} \\
\texttt{ZOMBIE} & \texttt{Z} & \texttt{Z} & Terminó pero el padre no hizo \texttt{wait()} \\
\bottomrule
\end{tabular}
\end{center}

El estado \texttt{TERMINATED} no aparece en \texttt{ps} porque, una vez
hecho el \texttt{wait()}, el descriptor se libera y el proceso desaparece
de la tabla. Los flags adicionales (\texttt{<}, \texttt{N}, \texttt{s},
\texttt{+}, \texttt{l}, \texttt{L}) refinan la información: prioridad
no estándar, líder de sesión, foreground, multi-hilo, páginas
\emph{locked} en RAM.

\subsection*{1.b Ordenar por CPU y memoria}

En Linux:

\begin{lstlisting}
ps aux --sort=-%cpu | head
ps aux --sort=-%mem | head
\end{lstlisting}

En macOS, \texttt{ps} hereda de BSD y no soporta \texttt{--sort}. El
equivalente es:

\begin{lstlisting}
ps aux -r | head -n 10   # por CPU descendente
ps aux -m | head -n 10   # por memoria residente descendente
\end{lstlisting}

\subsection*{1.c Procesos del usuario actual}

\begin{lstlisting}
ps -U "$USER" -o user,pid,%cpu,%mem,stat,command
\end{lstlisting}

Alternativamente, filtrando con \texttt{awk}:

\begin{lstlisting}
ps aux | awk -v u="$USER" '$1 == u'
\end{lstlisting}

\subsection*{Conexión con la teoría}

\texttt{ps} no es un comando autónomo: lee directamente las estructuras
del kernel que la teoría llama \emph{task descriptors}. En Linux la
información vive bajo \texttt{/proc/<PID>/} (un \emph{pseudo-filesystem}
expuesto por el kernel) y \texttt{ps} la formatea para el usuario. Cada
columna corresponde a un campo del descriptor: \texttt{PID}~$\rightarrow$
identificador, \texttt{STAT}~$\rightarrow$ campo \texttt{state},
\texttt{\%CPU}/\texttt{TIME}~$\rightarrow$ contadores acumulados por el
scheduler, \texttt{VSZ}/\texttt{RSS}~$\rightarrow$ \emph{mapa de memoria}
del proceso, \texttt{COMMAND}~$\rightarrow$ programa del que es instancia.

% =====================================================================
\section{Ejercicio 2 --- Control de un proceso con señales}
% =====================================================================

Se pide manipular el estado del script \texttt{runseconds.sh} (un
busy-wait de 120 segundos) lanzándolo en background, deteniéndolo con
\texttt{SIGSTOP}, reanudándolo con \texttt{SIGCONT} y finalmente
esperando o matándolo.

\subsection*{2.a El script}

\begin{lstlisting}[style=shellstyle]
#!/bin/bash
# Run for 120 seconds
end=$((SECONDS + 120))

while [ $SECONDS -lt $end ]; do
    : # this command does nothing
done

echo "Done!"
\end{lstlisting}

La variable \texttt{SECONDS} es interna del shell y avanza solo cuando
el shell se está ejecutando: si el proceso está en estado \texttt{T}
(detenido), \texttt{SECONDS} no avanza y la duración total se mide en
tiempo de CPU corrido, no en tiempo real.

\subsection*{2.b Lanzar en background y verificar \texttt{R}}

\begin{lstlisting}
chmod +x runseconds.sh
./runseconds.sh &
PID=$!
ps -o pid,stat,command -p $PID
\end{lstlisting}

\begin{lstlisting}[style=textstyle]
  PID STAT CMD
12345 R    bash ./runseconds.sh
\end{lstlisting}

\subsection*{2.c Pasar a \texttt{T} con \texttt{SIGSTOP}}

\begin{lstlisting}
kill -STOP $PID
ps -o pid,stat,command -p $PID
\end{lstlisting}

\begin{lstlisting}[style=textstyle]
  PID STAT COMMAND
12345 T    bash ./runseconds.sh
\end{lstlisting}

\texttt{SIGSTOP} (señal 19 en Linux) no puede ser ignorada ni capturada;
el kernel saca al proceso de toda cola de planificación hasta recibir
\texttt{SIGCONT}.

\subsection*{2.d Reanudar con \texttt{SIGCONT}}

\begin{lstlisting}
kill -CONT $PID
ps -o pid,stat,command -p $PID
\end{lstlisting}

El estado vuelve a \texttt{R}.

\subsection*{2.e Esperar o terminar}

\begin{lstlisting}
wait $PID         # esperar a que termine solo
kill $PID         # SIGTERM (terminacion ordenada)
kill -KILL $PID   # SIGKILL (no se puede ignorar)
\end{lstlisting}

\subsection*{Conexión con la teoría}

Cada paso del ejercicio corresponde a una transición concreta del
diagrama de estados (Figura 5 de las Notas 5--6):

\begin{center}
\small
\begin{tabular}{@{}llp{8.0cm}@{}}
\toprule
Paso & Señal & Transición en el modelo teórico \\
\midrule
2.b & (lanzamiento) & \texttt{NEW}~$\rightarrow$~\texttt{RUNNABLE}~$\rightarrow$~\texttt{RUNNING} \\
2.c & \texttt{SIGSTOP} & \texttt{RUNNING}~$\rightarrow$~\emph{Stopped} (fuera del scheduler) \\
2.d & \texttt{SIGCONT} & \emph{Stopped}~$\rightarrow$~\texttt{RUNNABLE}~$\rightarrow$~\texttt{RUNNING} \\
2.e & \texttt{SIGTERM}/fin & \texttt{RUNNING}~$\rightarrow$~\texttt{ZOMBIE}~$\rightarrow$~\texttt{TERMINATED} \\
\bottomrule
\end{tabular}
\end{center}

\texttt{SIGSTOP} es el ejemplo concreto del estado \emph{controlado por
el SO}: el proceso no espera I/O (no está \texttt{SLEEPING}) ni compite
por CPU (no está \texttt{RUNNABLE}); fue retirado de toda cola hasta
recibir \texttt{SIGCONT}. \texttt{SIGTERM} puede ser capturada por el
proceso para hacer una salida ordenada; \texttt{SIGKILL} y \texttt{SIGSTOP}
las maneja el kernel directamente y no pueden ser interceptadas, lo
que las convierte en herramientas seguras contra procesos colgados.

% =====================================================================
\section{Ejercicio 3 --- Árbol de procesos con \texttt{pstree}}
% =====================================================================

Se pide visualizar el árbol de procesos del sistema y determinar el
primer proceso lanzado y su PID.

\subsection*{3.a Salida de \texttt{pstree}}

En GNU/Linux \texttt{pstree} viene preinstalado como parte de
\emph{psmisc}. En macOS hace falta instalar la versión de Fred Hucht
con Homebrew (\texttt{brew install pstree}); las opciones difieren
ligeramente de las de Linux.

\begin{lstlisting}
pstree
\end{lstlisting}

Las primeras líneas son aproximadamente:

\begin{lstlisting}[style=textstyle]
-+= 00001 root /sbin/launchd        (macOS)
 |--= 00312 root /usr/libexec/logd
 |--= 00313 root /usr/libexec/smd
 ...
\end{lstlisting}

Los conectores \texttt{-+=} y \texttt{|--=} arman la jerarquía y los
números de cinco dígitos son los PIDs.

\subsection*{3.b El primer proceso}

El primer proceso de espacio de usuario tiene siempre \textbf{PID~1} y
es ancestro de todos los demás:

\begin{lstlisting}
ps -p 1 -o pid,ppid,user,comm
\end{lstlisting}

\begin{lstlisting}[style=textstyle]
  PID  PPID USER COMM
    1     0 root /sbin/launchd
\end{lstlisting}

El \texttt{PPID = 0} corresponde al proceso interno del kernel
(\emph{scheduler/swapper}) que no aparece en espacio de usuario. Quién
ejecuta como PID~1 depende del sistema: \texttt{launchd} en macOS,
\texttt{systemd} en Linux moderno, \texttt{init} (SysV/BusyBox) en
Linux tradicional, \texttt{init} en FreeBSD.

\subsection*{Conexión con la teoría}

La Figura 6 del capítulo \emph{Procesos y threads} describe la secuencia
de \emph{bootstrapping} en sistemas tipo UNIX:

\begin{lstlisting}[style=textstyle]
initcode  (pid=1, hardcoded en el kernel)
   | exec("init")
init      (pid=1)
   | fork(); child: exec("tty")
tty       (pid=2)
   | exec("login")
login     (pid=3)
   | exec("sh")
shell     (pid=4)
\end{lstlisting}

El kernel construye \texttt{initcode} en una página especial al terminar
el boot. Ese código invoca \texttt{exec("init")}, que reemplaza la imagen
del proceso PID~1 por el binario del \emph{init} del sistema. Toda otra
tarea de espacio de usuario surge por \texttt{fork()} desde algún
descendiente de PID~1: por construcción, \texttt{pstree} produce un
árbol y nunca un bosque.

Esta arquitectura justifica la \emph{adopción de huérfanos} por parte
de PID~1: si un proceso muriese sin hacer \texttt{wait()} sobre sus
hijos, estos quedarían en estado \texttt{ZOMBIE} indefinidamente.
Reasignando los huérfanos a init, el sistema asegura que tarde o
temprano alguien hará \texttt{wait()} y la tabla de procesos se mantenga
limpia.

% =====================================================================
\section{Ejercicio 4 --- Planificación con \texttt{nice} y \texttt{renice}}
% =====================================================================

Se pide ejecutar un proceso CPU-bound, lanzar una segunda instancia con
menor prioridad y observar las diferencias en \emph{context switches}
con \texttt{cat /proc/<PID>/sched | grep nr\_switches}. Después,
modificar la prioridad en caliente y, finalmente, lanzar un proceso
con prioridad superior.

\subsection*{4.a Programa de carga}

\begin{lstlisting}[style=shellstyle]
#!/bin/bash
# busyloop.sh: ciclo infinito, 100% CPU.
while :; do :; done
\end{lstlisting}

\subsection*{4.b Una instancia, observada con \texttt{top}}

\begin{lstlisting}
./busyloop.sh
# en otra terminal:
top   # tecla P ordena por %CPU; q sale
\end{lstlisting}

\texttt{bash busyloop.sh} consume cerca del 100\% de un núcleo, con
\texttt{NI=0} y \texttt{PR} por defecto.

\subsection*{4.c Segunda instancia con \texttt{nice +15}}

\begin{lstlisting}
nice -n 15 bash ./busyloop.sh
\end{lstlisting}

\texttt{top} muestra las dos instancias compitiendo: la nueva con
\texttt{NI=15} y \texttt{PR} numéricamente más alto (recordar que en
Linux un \texttt{PR} mayor significa \emph{menos} prioridad).

\paragraph{Cambios de contexto.} El kernel expone los context switches
acumulados de cada proceso en \texttt{/proc/<PID>/sched}:

\begin{lstlisting}
cat /proc/<PID>/sched | grep nr_switches
\end{lstlisting}

\textbf{¿El proceso de menor prioridad ejecuta más o menos cambios de
contexto por unidad de tiempo? Más.} El razonamiento usa la regla
central del CFS (\emph{Completely Fair Scheduler}) de Linux: el
\emph{quantum} asignado a una tarea es proporcional a su \emph{peso},
y el peso es decreciente con \texttt{nice}. Concretamente, el peso
aproximado es $1024 \cdot 1{.}25^{-\text{nice}}$, de modo que cada
incremento unitario de \texttt{nice} reduce el peso en $\sim 20\%$.
Un proceso con \texttt{nice=15} tiene un peso aproximadamente
$15$ veces menor que uno con \texttt{nice=0}, por lo que:

\begin{itemize}[leftmargin=*]
    \item recibe \emph{time slices} más cortos $\rightarrow$ es desalojado
        más seguido $\rightarrow$ \texttt{nr\_switches} crece más rápido
        por unidad de tiempo (paradójicamente, el \emph{menos prioritario}
        hace \emph{más} context switches);
    \item pero la cantidad \textbf{total} de CPU que recibe es menor.
\end{itemize}

Para verlo en vivo:

\begin{lstlisting}
watch -n1 "cat /proc/<PID_NICE15>/sched | grep nr_switches"
\end{lstlisting}

\subsection*{4.d Cambiar prioridad en caliente con \texttt{renice}}

\begin{lstlisting}
renice -n  5 -p <PID>     # subir nice (bajar prioridad)
renice -n -5 -p <PID>     # bajar nice (subir prioridad, requiere root)
ps -o pid,ni,pri,comm -p <PID>
\end{lstlisting}

\subsection*{4.e Mayor prioridad y por qué pide \texttt{sudo}}

\begin{lstlisting}
sudo nice -n -10 bash ./busyloop.sh
\end{lstlisting}

Bajar \texttt{nice} a un valor negativo aumenta la prioridad por encima
de la del resto de los procesos del usuario y, por eso, requiere
\textbf{root}. La razón es de \emph{seguridad y equidad}: si cualquier
usuario pudiera mejorarse a sí mismo la prioridad, podría desplazar a
procesos del sistema y monopolizar la CPU, transformando una limitación
del scheduler en una herramienta de \emph{denegación de servicio}. Por
eso la convención histórica de UNIX es que solo \texttt{root} puede
asignar \texttt{nice} negativos. Un usuario sin privilegios sí puede
\emph{subir} el \texttt{nice} de un proceso suyo (cederle prioridad a
otros), nunca lo opuesto. Esto refleja la regla general:
\emph{cada usuario puede ceder sus propios recursos pero no apropiarse
de los ajenos}.

\subsection*{Conexión con la teoría}

El \emph{quantum} es ``el período de tiempo máximo de uso de la CPU''
(Notas 5--6, \emph{Cambios de contexto}). La interrupción del
\emph{timer} es la implementación concreta del \emph{preemptive
multitasking}: gracias a ella el kernel puede recuperar la CPU sin la
cooperación del proceso, evitando que un thread que no invoca
\texttt{yield()} voluntariamente monopolice el procesador. Cada
desalojo es un \emph{context switch involuntario} contabilizado en
\texttt{nr\_switches}.

\paragraph{Limpieza.}

\begin{lstlisting}
pkill -f busyloop.sh
\end{lstlisting}

% =====================================================================
\section{Ejercicio 5 --- \texttt{counter.c}, race condition y \texttt{flock}}
% =====================================================================

Se pide compilar \texttt{counter.c}, observar que ejecuciones secuenciales
producen 2000 pero ejecuciones concurrentes pierden incrementos, explicar
la causa y solucionarlo con \texttt{flock()}. Por último, justificar si
\texttt{fd} es una variable compartida entre las distintas instancias.

\subsection*{5.a El programa}

\begin{lstlisting}[style=cstyle]
#include <fcntl.h>
#include <stdlib.h>
#include <string.h>
#include <unistd.h>
#include <stdio.h>
#include <sys/file.h>

#define N 10

int main()
{
    int pid = getpid(), i;
    int fd = open("counter.dat", O_RDWR);
    (void)pid;

    for (i=0; i<1000; i++) {
        char s[N];
        int n;

        lseek(fd, 0, SEEK_SET);
        read(fd, s, N);
        n = atoi(s);
        sprintf(s, "%d", n+1);
        lseek(fd, 0, SEEK_SET);   // rewind
        write(fd, s, strlen(s));
        fsync(fd);
    }
    close(fd);
    return 0;
}
\end{lstlisting}

\begin{lstlisting}
gcc -Wall -Wextra -o counter counter.c
echo 0 > counter.dat
\end{lstlisting}

\subsection*{5.b Ejecución secuencial}

\begin{lstlisting}
echo 0 > counter.dat
./counter ; ./counter
cat counter.dat        # 2000
\end{lstlisting}

Cada ejecución empieza después de que la anterior cerró el archivo, así
que ningún incremento se pierde.

\subsection*{5.c Ejecución concurrente sin lock}

\begin{lstlisting}
echo 0 > counter.dat
./counter & ./counter
wait
cat counter.dat        # < 2000, no determinista
\end{lstlisting}

El resultado es menor a $2000$ y cambia entre corridas: $1826$, $1740$,
$1766$, etc.

\paragraph{Por qué se pierden incrementos.} La secuencia
\emph{leer-modificar-escribir} no es atómica. Con dos procesos $A$ y $B$
en paralelo, una traza típica es:

\begin{lstlisting}[style=textstyle]
A lseek(0); read -> "100"          (v=100)
                            B lseek(0); read -> "100"  (v=100)
A sprintf("101"); write 3 bytes
                            B sprintf("101"); write 3 bytes
                                <-- pisa con el mismo valor
\end{lstlisting}

Es el patrón clásico de \emph{lost update}: ambos procesos leen el mismo
valor, ambos escriben el mismo resultado, y un incremento se pierde
silenciosamente. \texttt{fsync()} solo garantiza persistencia en disco;
no excluye a otro proceso que esté queriendo escribir.

\subsection*{5.d Solución con \texttt{flock()}}

\begin{lstlisting}[style=cstyle]
for (i=0; i<1000; i++) {
    char s[N];
    int n;

    flock(fd, LOCK_EX);

    lseek(fd, 0, SEEK_SET);
    read(fd, s, N);
    n = atoi(s);
    sprintf(s, "%d", n+1);
    lseek(fd, 0, SEEK_SET);
    write(fd, s, strlen(s));
    fsync(fd);

    flock(fd, LOCK_UN);
}
\end{lstlisting}

\begin{lstlisting}
echo 0 > counter.dat
./counter_locked & ./counter_locked
wait
cat counter.dat        # 2000, deterministico
\end{lstlisting}

\texttt{flock(fd, LOCK\_EX)} toma un lock exclusivo \emph{advisory} sobre
el descriptor. Mientras un proceso lo retiene, cualquier otro que intente
\texttt{LOCK\_EX} se duerme hasta que el primero llame \texttt{LOCK\_UN}.
La región crítica queda lo más corta posible: solo el read-modify-write,
no todo el bucle, lo que preserva el paralelismo cuando la región no se
está ejecutando.

\subsection*{5.e ¿Es \texttt{fd} una variable compartida?}

\textbf{No.} Cuando dos procesos hacen \texttt{./counter} por separado,
cada uno corre en su propio espacio de direcciones y cada uno hace su
propio \texttt{open("counter.dat", O\_RDWR)}. El kernel les devuelve un
descriptor en la \emph{tabla de descriptores} \emph{de cada proceso},
que casualmente probablemente sea \texttt{3} en ambos (porque
\texttt{0}, \texttt{1}, \texttt{2} ya están ocupados por
\texttt{stdin/stdout/stderr}), pero son entradas en tablas distintas.
Apuntan al mismo \emph{inodo} del filesystem, no a la misma
\texttt{struct file} del kernel; cada uno tiene su propio offset.

Esto explica por qué la race se da: como los offsets son independientes,
ambos procesos pueden hacer \texttt{lseek(0)} y \texttt{read} sin que el
kernel los serialice. El recurso compartido entre las dos instancias
\textbf{no es \texttt{fd}}, sino el \textbf{archivo en disco}, y la
sincronización tiene que hacerse sobre ese archivo (con \texttt{flock}).

\subsection*{Conexión con la teoría}

Las Notas 5--6 (\emph{Concurrencia e IPC}) definen la \emph{condición
de carrera} como ``el comportamiento del sistema que depende de la
secuencia de la ocurrencia de eventos externos que pueden generar flujos
de ejecución no deseados''. En este ejercicio la \emph{región crítica}
es el bloque \texttt{lseek + read + atoi + sprintf + lseek + write}, el
\emph{recurso compartido} es \texttt{counter.dat}, y el \emph{interleaving
no deseado} es el \emph{lost update} mostrado arriba.

\texttt{flock()} es la implementación concreta del patrón
\emph{acquire-uso-release} (Listado 1 del capítulo \emph{Locks}):

\begin{lstlisting}[style=cstyle]
acquire(&l);     // flock(fd, LOCK_EX)
use_resource(r); // operacion sobre el archivo
release(&l);    // flock(fd, LOCK_UN)
\end{lstlisting}

provee \emph{exclusión mutua} sobre la región crítica.

% =====================================================================
\section{Ejercicio 6 --- \texttt{pthreads-example.c}, mutex y stacks}
% =====================================================================

Se pide ejecutar el programa, observar que falla, corregirlo con
\texttt{pthread\_mutex\_t}, mostrar la dirección de la variable local
\texttt{tid} en cada hilo y verificar experimentalmente que todos los
hilos ejecutan en el mismo espacio de memoria del proceso.

\subsection*{6.a Programa original y bug}

\begin{lstlisting}[style=cstyle]
int counter = 0; // Shared resource

void* increment(void* thread_id) {
    int tid = *((int *) thread_id);
    for (int i = 0; i < 100000; i++) {
        if (i == 0) {
            printf("Thread %d running...\n", tid);
            printf("Address of local variable: %p\n", &tid);
        }
        counter++;
    }
    return NULL;
}

int main() {
    pthread_t t1, t2;
    int tn1 = 1, tn2 = 2;

    pthread_create(&t1, NULL, increment, &tn1);
    pthread_create(&t2, NULL, increment, &tn2);

    pthread_join(t1, NULL);
    pthread_join(t2, NULL);

    printf("Final Counter Value: %d\n", counter);
    return 0;
}
\end{lstlisting}

\begin{lstlisting}
gcc -Wall -Wextra -pthread -o pthreads-example pthreads-example.c
./pthreads-example
\end{lstlisting}

Salida (ejemplo):

\begin{lstlisting}[style=textstyle]
Thread 1 running...
Address of local variable: 0x16fcdefb4
Thread 2 running...
Address of local variable: 0x16fd6afb4
Final Counter Value: 102483
\end{lstlisting}

El valor final es menor a $200000$. La causa es que \texttt{counter++}
no es atómico: en máquina se traduce aproximadamente a tres instrucciones
\emph{load-modify-store}:

\begin{lstlisting}[style=textstyle]
lw   t0, counter
addi t0, t0, 1
sw   t0, counter
\end{lstlisting}

Dos hilos pueden leer el mismo valor antes de que cualquiera escriba; el
último \texttt{store} pisa al otro. Como ambos hilos viven en el mismo
proceso, comparten la misma página de memoria que aloja \texttt{counter}:
no hay copias separadas, hay una única palabra de memoria.

\subsection*{6.b Corrección con \texttt{pthread\_mutex\_t}}

\begin{lstlisting}[style=cstyle]
pthread_mutex_t mtx = PTHREAD_MUTEX_INITIALIZER;

for (int i = 0; i < 100000; i++) {
    if (i == 0) {
        printf("Thread %d running...\n", tid);
        printf("Address of local variable: %p\n", &tid);
    }
    pthread_mutex_lock(&mtx);
    counter++;
    pthread_mutex_unlock(&mtx);
}
\end{lstlisting}

\begin{lstlisting}[style=textstyle]
Thread 1 running...
Address of local variable: 0x16f702fb4
Thread 2 running...
Address of local variable: 0x16f78efb4
Final Counter Value: 200000
\end{lstlisting}

Siempre $200000$, de manera determinista. El mutex implementa un
\emph{sleep lock}: si un hilo intenta tomar un lock ocupado, el kernel
lo pone en \texttt{SLEEPING} hasta que el otro hilo lo libere. Esto
contrasta con un \emph{spinlock}, que haría busy-wait.

\subsection*{6.c Direcciones de \texttt{tid} y distancia entre stacks}

\begin{lstlisting}[style=textstyle]
Thread 2 running...
Address of local variable: 0x16fb26fb4
Thread 1 running...
Address of local variable: 0x16fa9afb4
\end{lstlisting}

\begin{lstlisting}[style=textstyle]
|&tid_2 - &tid_1| = |0x16fb26fb4 - 0x16fa9afb4| = 0x8c000
                  ~ 573 440 bytes ~ 560 KB
\end{lstlisting}

El valor exacto varía entre ejecuciones (ASLR del kernel), pero la
\emph{distancia} entre stacks consecutivos es constante y corresponde
al \textbf{tamaño del stack reservado por hilo}:

\begin{itemize}[leftmargin=*]
    \item macOS (Darwin): $\sim 512$ KB para hilos no-main.
    \item Linux/glibc: $8$ MB ($0\text{x}800000$, equivalente a
        \texttt{ulimit -s = 8192} KB).
    \item FreeBSD: $2$ MB.
\end{itemize}

\textbf{¿Están en la misma pila? No.} Cada hilo tiene su propio stack,
asignado por la biblioteca de \emph{pthreads} con \texttt{mmap} privado
al momento de \texttt{pthread\_create}. La distancia observada confirma
que son stacks separados.

\subsection*{6.d Experimento: todos los hilos comparten el address space}

Para verlo de manera directa, basta con que cada hilo imprima la
dirección del global \texttt{counter}, su PID, y la dirección de su
\texttt{tid} local:

\begin{lstlisting}[style=cstyle]
printf("[T%d] pid=%d, &counter=%p, &tid=%p\n",
       tid, getpid(), (void*)&counter, (void*)&tid);
\end{lstlisting}

\begin{lstlisting}[style=textstyle]
[T1] pid=44213, &counter=0x10054af80, &tid=0x16fa9afb4
[T2] pid=44213, &counter=0x10054af80, &tid=0x16fb26fb4
\end{lstlisting}

Lo observado:

\begin{itemize}[leftmargin=*]
    \item \textbf{Mismo PID} $\rightarrow$ son tareas dentro del mismo
        proceso.
    \item \textbf{Misma dirección para \texttt{\&counter}} $\rightarrow$
        la variable global vive en la misma página virtual; ambos hilos
        usan la misma tabla de páginas. \emph{Por eso} el
        \texttt{counter++} sin lock genera la race condition.
    \item \textbf{\texttt{\&tid} distintos} $\rightarrow$ stacks
        independientes.
\end{itemize}

\subsection*{Conexión con la teoría}

El experimento es la confirmación empírica del modelo de proceso
multithreaded de las Notas 5--6 (Figura 1):
``cada thread usa su propia pila aunque todos los threads comparten el
espacio de memoria del proceso''. Texto, datos, heap y variables
globales como \texttt{counter} son \textbf{compartidos} entre todos los
hilos. El stack es \textbf{privado} a cada hilo. La elección entre
spinlock y sleep lock (mutex) sigue la regla de la teoría: spinlock
para regiones cortas, mutex para regiones largas o con alta contención.

% =====================================================================
\section{Ejercicio 7 --- Multithreading en Java}
% =====================================================================

Se pide analizar (1) los mecanismos de creación y control de threads en
Java, (2) los mecanismos de sincronización y su clasificación, y (3)
implementar el equivalente del Ejercicio~6 sin race conditions.

\subsection*{7.a Creación y control de threads}

Java provee soporte de hilos de \emph{primera clase} en la JVM. Cada
hilo Java se corresponde uno-a-uno con un hilo nativo del sistema
operativo. Hay dos formas equivalentes de crearlos:

\begin{lstlisting}[style=javastyle]
// Forma 1: extender Thread (poco flexible)
class MiHilo extends Thread {
    @Override public void run() { /* ... */ }
}
new MiHilo().start();

// Forma 2: implementar Runnable (preferida)
class Tarea implements Runnable {
    @Override public void run() { /* ... */ }
}
new Thread(new Tarea()).start();
new Thread(() -> { /* ... */ }).start();   // con lambda
\end{lstlisting}

Métodos clave:

\begin{center}
\small
\begin{tabular}{@{}lp{8.5cm}@{}}
\toprule
Método & Qué hace \\
\midrule
\texttt{start()}        & Pide a la JVM crear el hilo y ejecutar \texttt{run()} en paralelo \\
\texttt{run()}          & Cuerpo del hilo. Llamarlo directo NO crea hilo \\
\texttt{join()}         & Bloquea al llamador hasta que el hilo termine \\
\texttt{Thread.sleep(ms)} & Pone al hilo actual en \texttt{TIMED\_WAITING} \\
\texttt{Thread.yield()} & Sugiere ceder la CPU \\
\texttt{interrupt()}    & Marca el hilo como interrumpido (lo despierta de \texttt{sleep}/\texttt{wait}) \\
\texttt{setDaemon(true)} & El hilo no impide la terminación del proceso \\
\bottomrule
\end{tabular}
\end{center}

El ciclo de vida de un hilo Java tiene los estados \texttt{NEW},
\texttt{RUNNABLE}, \texttt{BLOCKED}, \texttt{WAITING},
\texttt{TIMED\_WAITING} y \texttt{TERMINATED}, alineado con el modelo
de la teoría (la JVM agrupa ``ejecutando'' y ``listo'' en
\texttt{RUNNABLE}).

\subsection*{7.b Sincronización: \texttt{synchronized} y \texttt{wait}/\texttt{notify}}

Java provee dos primitivas integradas, sostenidas por una abstracción
común: el \emph{monitor}. Cada objeto Java lleva implícitamente un
monitor (mutex + variable de condición) que se activa con la palabra
clave \texttt{synchronized}.

\begin{lstlisting}[style=javastyle]
// Bloque sincronizado: el lock es el objeto entre parentesis.
synchronized (obj) { /* region critica */ }

// Metodo de instancia sincronizado: lock implicito = this.
public synchronized void m() { /* ... */ }

// Metodo estatico sincronizado: lock implicito = Class del tipo.
public static synchronized void m() { /* ... */ }
\end{lstlisting}

Los locks son \emph{reentrantes} y se liberan automáticamente al salir
del bloque, incluso por excepción.

\paragraph{\texttt{wait}/\texttt{notify}/\texttt{notifyAll}.} Solo
pueden invocarse mientras se posee el monitor del objeto.
\texttt{obj.wait()} libera atómicamente el monitor de \texttt{obj} y
duerme al hilo. \texttt{obj.notify()} despierta a uno (cualquiera) de
los hilos esperando; \texttt{obj.notifyAll()}, a todos.

La regla canónica es esperar siempre dentro de un \texttt{while} (no
de un \texttt{if}) por \emph{spurious wakeups} y por el patrón
multi-productor/multi-consumidor:

\begin{lstlisting}[style=javastyle]
synchronized (obj) {
    while (!condicion) obj.wait();
    // ahora la condicion es cierta y tengo el monitor
}
\end{lstlisting}

\paragraph{Clasificación.} Las primitivas se agrupan en dos familias
clásicas: exclusión mutua y sincronización por condición. Java cubre
ambas:

\begin{center}
\small
\begin{tabular}{@{}lll@{}}
\toprule
Mecanismo Java & Categoría & Equivalente POSIX \\
\midrule
\texttt{synchronized} & Exclusión mutua (mutex/monitor) & \texttt{pthread\_mutex\_t} / \texttt{flock} \\
\texttt{wait}/\texttt{notify} & Variable de condición & \texttt{pthread\_cond\_wait}/\texttt{cond\_signal} \\
\texttt{Thread.join()} & Sincronización por finalización & \texttt{pthread\_join} / \texttt{wait()} syscall \\
\texttt{volatile} & Visibilidad entre threads & \texttt{\_\_atomic\_*} / barreras de memoria \\
\bottomrule
\end{tabular}
\end{center}

Además, \texttt{java.util.concurrent} ofrece \texttt{ReentrantLock},
\texttt{Semaphore}, \texttt{CountDownLatch}, \texttt{CyclicBarrier},
\texttt{BlockingQueue}, etc.

\subsection*{7.c Equivalente Java sin race conditions}

\begin{lstlisting}[style=javastyle]
public class CounterExample {
    private static int counter = 0;
    private static final Object lock = new Object();
    private static final int N = 100_000;

    public static void main(String[] args) throws InterruptedException {
        Thread t1 = new Thread(new Incrementer(1));
        Thread t2 = new Thread(new Incrementer(2));

        t1.start();
        t2.start();

        t1.join();
        t2.join();

        System.out.println("Final Counter Value: " + counter);
    }

    static class Incrementer implements Runnable {
        private final int tid;

        Incrementer(int tid) { this.tid = tid; }

        @Override
        public void run() {
            System.out.println("Thread " + tid + " running...");
            for (int i = 0; i < N; i++) {
                synchronized (lock) {
                    counter++;
                }
            }
        }
    }
}
\end{lstlisting}

\begin{lstlisting}
javac CounterExample.java
java CounterExample
\end{lstlisting}

\begin{lstlisting}[style=textstyle]
Thread 1 running...
Thread 2 running...
Final Counter Value: 200000
\end{lstlisting}

Por qué funciona: \texttt{synchronized (lock)} garantiza
\textbf{exclusión mutua} (un solo hilo a la vez en el bloque) y
\textbf{visibilidad} (el modelo de memoria de Java establece la relación
\emph{happens-before} entre \texttt{unlock} y el siguiente \texttt{lock}).
El incremento se vuelve atómico respecto de los otros hilos y el
resultado es determinístico.

\subsection*{Conexión con la teoría}

Los bloques \texttt{synchronized} de Java son la implementación a nivel
de lenguaje de los \emph{monitores} descritos por la teoría
(Notas 5--6, \emph{Monitores}): cada objeto Java es un monitor con
mutex implícito y una variable de condición asociada
(\texttt{wait}/\texttt{notify}). La equivalencia conceptual es exacta:
\texttt{synchronized (obj)} corresponde a \texttt{acquire(\&obj.lock)};
salir del bloque, a \texttt{release(\&obj.lock)};
\texttt{obj.wait()}, a \texttt{wait(\&obj.cond, \&obj.lock)};
\texttt{obj.notify()}, a \texttt{signal(\&obj.cond)}.

% =====================================================================
\section{Ejercicio 8 --- Operaciones del kernel ante interrupción del timer}
% =====================================================================

Se pide describir las operaciones del kernel cuando un proceso es
interrumpido por el \emph{timer} y ya consumió su \emph{quantum},
hasta retornar de la interrupción, indicando en qué estado el kernel
deja al proceso. La descripción sigue el modelo xv6 sobre RISC-V.

\subsection*{Punto de partida}

Un proceso $P$ está corriendo en modo usuario cuando ocurre una IRQ
del timer (CLINT). La IRQ está delegada de M-mode a S-mode y
configurada por el kernel al inicializarse.

\subsection*{Paso 1. La CPU atrapa la interrupción (hardware)}

\begin{enumerate}[leftmargin=*,nosep]
    \item Verifica que las interrupciones del modo S estén habilitadas
        (\texttt{sstatus.SIE = 1}).
    \item Salva en \texttt{sepc} la próxima instrucción de $P$.
    \item Salva el modo previo en \texttt{sstatus.SPP} (\texttt{U}).
    \item Codifica el motivo en \texttt{scause} (timer en supervisor).
    \item Cambia el modo a supervisor.
    \item Salta a \texttt{stvec} (que apuntaba a \texttt{uservec}).
\end{enumerate}

\subsection*{Paso 2. \texttt{uservec}}

Salva todos los registros de propósito general en el \texttt{trapframe}
del proceso, conmuta a la pagetable del kernel, carga el stack del
kernel y salta a \texttt{usertrap()}.

\subsection*{Paso 3. \texttt{usertrap()}}

Setea \texttt{stvec = kernelvec}, salva \texttt{sepc} en el trapframe,
identifica la causa con \texttt{devintr()} y, al ser timer, llama a
\texttt{yield()}.

\subsection*{Paso 4. \texttt{yield()} cambia el estado a \texttt{RUNNABLE}}

\begin{enumerate}[leftmargin=*,nosep]
    \item Toma \texttt{p->lock}.
    \item Cambia \texttt{p->state} de \texttt{RUNNING} a
        \textbf{\texttt{RUNNABLE}}.
    \item Llama \texttt{sched()}, que termina haciendo
        \texttt{swtch(\&p->context, \&cpu->scheduler)}.
\end{enumerate}

A partir de aquí, $P$ está \textbf{congelado en \texttt{RUNNABLE}}, sin
CPU.

\subsection*{Pasos 5--6. El scheduler ejecuta otros procesos y luego vuelve}

\texttt{scheduler()} elige una tarea \texttt{RUNNABLE} $Q$ y le da CPU.
Más adelante el scheduler vuelve a seleccionar $P$ y hace
\texttt{swtch(\&cpu->scheduler, \&p->context)}, reanudándolo dentro de
\texttt{sched()}.

\subsection*{Pasos 7--8. \texttt{usertrapret} $\rightarrow$ \texttt{userret} $\rightarrow$ \texttt{sret}}

\texttt{usertrapret()} prepara el regreso (\texttt{stvec},
\texttt{sepc}, \texttt{sstatus.SPP=0}, configura el trapframe).
\texttt{userret} conmuta a la pagetable del proceso, restaura los
registros y ejecuta \texttt{sret}, que en una sola operación atómica
restaura \texttt{pc} desde \texttt{sepc}, restaura el modo a \texttt{U}
y rehabilita interrupciones. $P$ vuelve a ejecutar como si nada
hubiese ocurrido.

\subsection*{Estado final del proceso}

A lo largo de la rutina hay dos estados clave:

\begin{itemize}[leftmargin=*,nosep]
    \item Mientras $P$ está fuera de la CPU (entre los pasos 4 y 6), su
        estado en el PCB es \textbf{\texttt{RUNNABLE}}. Esa es la
        respuesta directa al ejercicio: el kernel deja a $P$ en
        \texttt{RUNNABLE} después de consumir su quantum.
    \item Cuando finalmente retorna de la interrupción, $P$ está
        nuevamente en \texttt{RUNNING}.
\end{itemize}

La interrupción del timer no \emph{bloquea} al proceso: solo lo
desaloja para darle CPU a otro. Por eso el estado correcto es
\texttt{RUNNABLE} (sigue siendo elegible para ser planificado), no
\texttt{SLEEPING} (que se reservaría para esperas por eventos como I/O,
caso del Ejercicio~9).

\subsection*{Ruta completa, esquemática}

\begin{lstlisting}[style=textstyle]
[user mode, P]
   | HW: timer irq -> sepc, SPP, scause, salta a stvec
   v
uservec (trampoline.S)        <-- salva regs en trapframe
   |
usertrap() en trap.c           <-- reconoce timer via devintr()
   |
yield() -> p->state = RUNNABLE  <-- aca el kernel deja a P
   |
sched() -> swtch(&p->context, &cpu->scheduler)   <-- 1er ctx switch
   v
scheduler() elige otro Q
   | swtch(&cpu->scheduler, &Q->context)
   |  ... eventualmente, scheduler vuelve a P ...
   v
swtch() reanuda en sched()/yield() de P
   v
usertrapret -> userret -> sret  <-- vuelve a modo usuario
   v
[user mode, P sigue ejecutando]
\end{lstlisting}

% =====================================================================
\section{Ejercicio 9 --- Operaciones del kernel ante \texttt{read} sobre disco}
% =====================================================================

Se pide la misma descripción del Ejercicio~8 pero cuando el proceso
ejecuta el \emph{syscall} \texttt{read} sobre un archivo en disco
(operación lenta) y describir cómo continúa la ejecución.

\subsection*{Punto de partida}

\begin{lstlisting}[style=cstyle]
n = read(fd, buf, len);
\end{lstlisting}

La función de la libc deja los argumentos en registros y dispara
\texttt{ECALL} (RISC-V) para entrar al kernel.

\subsection*{Paso 1. Trap por syscall (HW + uservec + usertrap)}

Identico al Ejercicio~8 pero \texttt{scause = 8} (\emph{environment
call from U-mode}). \texttt{usertrap()} reconoce que es syscall, avanza
\texttt{sepc} para que la vuelta apunte a la instrucción siguiente al
\texttt{ECALL}, habilita interrupciones y llama a \texttt{syscall()}.

\subsection*{Paso 2. Despacho a \texttt{sys\_read}}

\texttt{syscall()} consulta \texttt{a7} y ejecuta \texttt{sys\_read()},
que recupera \texttt{fd}, \texttt{buf}, \texttt{len} del trapframe,
valida y llama a \texttt{fileread()} $\rightarrow$ \texttt{readi()}
$\rightarrow$ \texttt{bread(dev, blockno)}.

\subsection*{Paso 3. Cache miss: arranca I/O y se duerme}

\texttt{bread()} busca el bloque en la \emph{buffer cache}.

\begin{itemize}[leftmargin=*,nosep]
    \item Cache hit: el bloque está en RAM, se sigue copiando datos al
        usuario sin bloquearse.
    \item Cache miss: \texttt{virtio\_disk\_rw(bp, 0)} encola la
        lectura al controlador VirtIO y luego:
\end{itemize}

\begin{lstlisting}[style=cstyle]
while ((bp->flags & B_VALID) == 0)
    sleep(bp, &lock);
\end{lstlisting}

\texttt{sleep(chan, lock)} es la primitiva clave:

\begin{enumerate}[leftmargin=*,nosep]
    \item Toma \texttt{p->lock}.
    \item Libera el lock que protege la condición.
    \item Setea \texttt{p->chan = bp} (el ``canal'' sobre el que duerme).
    \item Setea \textbf{\texttt{p->state = SLEEPING}}.
    \item Llama \texttt{sched()} $\rightarrow$ \texttt{swtch}.
\end{enumerate}

A partir de aquí $P$ está \texttt{SLEEPING} en una \emph{wait queue}.
Es justamente esto lo que distingue al \texttt{read} con I/O del simple
agotamiento de quantum: como la operación puede tardar mucho, el kernel
\textbf{no} deja al proceso \texttt{RUNNABLE}; lo bloquea y le cede la
CPU a quien la pueda aprovechar.

\subsection*{Paso 4. El scheduler corre otros procesos}

\texttt{scheduler()} elige un \texttt{RUNNABLE} y le da CPU.
Mientras tanto, el disco procesa la lectura \emph{en paralelo} a la CPU.

\subsection*{Paso 5. El disco termina y manda IRQ}

\texttt{devintr()} reconoce la IRQ del disco y llama a
\texttt{virtio\_disk\_intr()}, que marca el buffer como
\texttt{B\_VALID} y ejecuta \texttt{wakeup(bp)}: para todo proceso con
\texttt{state == SLEEPING \&\& chan == bp}, le cambia el estado a
\textbf{\texttt{RUNNABLE}}.

\subsection*{Paso 6. El scheduler vuelve a $P$}

\texttt{swtch(\&cpu->scheduler, \&p->context)} reanuda $P$ dentro de
\texttt{sleep()}. \texttt{sleep} retoma el lock, marca
\texttt{p->state = RUNNING} y retorna al \texttt{while} de
\texttt{bread()}. Como ahora \texttt{B\_VALID} está seteado, el
\texttt{while} termina, \texttt{bread} retorna el buffer y \texttt{readi}
copia los bytes al \texttt{buf} de usuario con \texttt{copyout()}.

\subsection*{Paso 7. Retorno a modo usuario}

El valor de retorno (bytes leídos) se escribe en
\texttt{p->trapframe->a0} y se retorna por la ruta
\texttt{usertrap}~$\rightarrow$~\texttt{usertrapret}~$\rightarrow$~\texttt{userret}~$\rightarrow$~\texttt{sret}.

\subsection*{Estados por los que pasó $P$}

\begin{center}
\small
\begin{tabular}{@{}ll@{}}
\toprule
Etapa & Estado de $P$ \\
\midrule
Antes del \texttt{read}                          & \texttt{RUNNING} \\
Dentro de \texttt{sys\_read}/\texttt{bread}, esperando I/O & \textbf{\texttt{SLEEPING}} \\
Después del \texttt{wakeup(bp)}                  & \textbf{\texttt{RUNNABLE}} \\
Cuando el scheduler lo vuelve a elegir            & \texttt{RUNNING} \\
\bottomrule
\end{tabular}
\end{center}

\subsection*{Conexión con la teoría}

La traza expone el complemento exacto del Ejercicio~8 y reproduce el
esquema teórico de las primitivas \texttt{suspend}/\texttt{wakeup}
(Notas 5--6, \emph{Implementación de tareas en el kernel}):

\begin{lstlisting}[style=cstyle]
// suspend current task in wait queue and yield the cpu
void suspend(struct wait_queue *wq) {
    current->state = WAITING;
    enqueue(wq, current);
    yield();
}

// wakeup a process waiting in wait queue
void wakeup(struct wait_queue *wq) {
    struct task *task = dequeue(wq);
    if (task) task->state = RUNNABLE;
}
\end{lstlisting}

La diferencia clave con el ejercicio anterior: el proceso no es
desalojado por \emph{quantum}; se bloquea voluntariamente porque no
puede progresar hasta que llegue el dato. Por eso la transición es a
\texttt{SLEEPING} y no a \texttt{RUNNABLE}. La cooperación
kernel + dispositivo + scheduler es lo que permite la concurrencia
entre I/O y CPU.

% =====================================================================
\section{Ejercicio 10 --- Demostración semi-formal de Peterson}
% =====================================================================

Se pide mostrar de manera semi-formal que la solución de Peterson
(con dos procesos) garantiza \emph{exclusión mutua} y \emph{progreso}.

\subsection*{El algoritmo}

\begin{lstlisting}[style=cstyle]
// Variables compartidas, inicializadas en falso / 0
bool flag[2] = { false, false };
int  turn   = 0;          // o 1, da igual

// Codigo que ejecuta el proceso i (j = 1 - i es el otro)
void p_i(void) {
    while (true) {
        flag[i] = true;            // (1) anuncio interes
        turn    = j;               // (2) cedo el turno al otro
        while (flag[j] && turn == j)   // (3) busy wait
            ;
        // ----- seccion critica -----
        flag[i] = false;           // (4) salgo de la SC
        // ----- seccion no critica -----
    }
}
\end{lstlisting}

\paragraph{Hipótesis del modelo.}
\begin{itemize}[leftmargin=*,nosep]
    \item Las lecturas y escrituras a \texttt{flag[i]} y a \texttt{turn}
        son atómicas.
    \item Solo $P_i$ escribe en \texttt{flag[i]}; \texttt{turn} la
        escriben ambos.
    \item No se asume orden estricto de ejecución entre $P_0$ y $P_1$:
        el scheduler puede intercalar arbitrariamente.
\end{itemize}

\subsection*{Parte A. Exclusión mutua}

\textbf{Afirmación.} En ningún estado alcanzable del sistema, ambos
procesos pueden estar simultáneamente dentro de la sección crítica.

\textbf{Demostración por contradicción.} Supongamos que en algún
instante $t$, $P_0$ y $P_1$ están ambos en la SC. Justo antes de
entrar:

\begin{itemize}[leftmargin=*,nosep]
    \item $P_0$ salió de su \texttt{while} (3), por lo que la guarda
        falló:
        \[
        \neg(\text{flag}[1] = \text{true} \land \text{turn} = 1)
        \;\Rightarrow\;
        \text{flag}[1] = \text{false} \;\lor\; \text{turn} = 0 \quad (\star)
        \]
    \item Simétricamente para $P_1$:
        \[
        \text{flag}[0] = \text{false} \;\lor\; \text{turn} = 1
        \quad (\star\star)
        \]
\end{itemize}

Como ambos están \emph{dentro} de la SC, ambos ya ejecutaron su paso
(1):
\[
\text{flag}[0] = \text{true} \quad \text{y} \quad \text{flag}[1] = \text{true}.
\]
(Ninguno limpia su flag hasta el paso (4), que está después de la SC.)

Reemplazando en $(\star)$ y $(\star\star)$:

\begin{itemize}[leftmargin=*,nosep]
    \item De $(\star)$: como \texttt{flag[1] = true}, debe ser
        \texttt{turn = 0}.
    \item De $(\star\star)$: como \texttt{flag[0] = true}, debe ser
        \texttt{turn = 1}.
\end{itemize}

Por lo tanto \texttt{turn = 0} y \texttt{turn = 1} simultáneamente.
Contradicción. $\blacksquare$

\paragraph{¿Por qué el orden (1) $\rightarrow$ (2) es esencial?} Si se
invirtieran (primero \texttt{turn = j}, después \texttt{flag[i] = true}),
la prueba se rompe: un proceso podría leer la flag del otro antes de
que el otro haya anunciado interés, y ambos podrían colarse
simultáneamente.

\subsection*{Parte B. Progreso}

\textbf{Afirmación.} Si la SC está libre y al menos un proceso quiere
entrar, alguno entrará en tiempo finito. El único proceso que puede
impedirle el ingreso a otro es uno que también quiere entrar.

\paragraph{Caso 1: $P_0$ quiere entrar y $P_1$ está en la sección no
crítica.} $P_1$ no está en (1)--(3): por lo tanto $\text{flag}[1] =
\text{false}$. La guarda del \texttt{while} de $P_0$ evalúa
$\text{flag}[1] \land \text{turn} = 1 \;\Rightarrow\;
\text{false}$. $P_0$ entra. $P_1$ no influye.

\paragraph{Caso 2: ambos procesos quieren entrar.} Ambos ejecutaron (1),
ambos \texttt{flag} quedaron en \texttt{true}. Sea cual sea el valor
final de \texttt{turn} antes de que alguno entre al \texttt{while}:

\begin{itemize}[leftmargin=*,nosep]
    \item Si \texttt{turn = 1}: la guarda de $P_0$ es verdadera (gira),
        la de $P_1$ es falsa. \textbf{$P_1$ entra.}
    \item Si \texttt{turn = 0}: simétricamente, \textbf{$P_0$ entra.}
\end{itemize}

En cualquier caso hay un ganador. Como \texttt{turn} ya no se vuelve a
tocar mientras quien tiene el turno esté esperando o dentro de la SC,
el ganador no es derrotado por nuevas escrituras.

\paragraph{Caso 3: el proceso que estaba en la SC ya terminó.} Al salir
ejecutó (4) (\texttt{flag[i] = false}); la guarda del otro ve
\texttt{flag[i] = false}, falla, y entra. Ningún proceso queda atrapado
por uno que ya terminó la SC. $\blacksquare$

\subsection*{Espera acotada (bonus)}

Para dos procesos la cota es trivial: a lo sumo el otro proceso entra
\emph{una sola vez} entre que yo expreso interés y entro yo. Cuando
$P_j$ sale de la SC y vuelve a competir, ejecuta nuevamente (2) con
\texttt{turn = i}; a partir de ahí mi \texttt{while} falla y entro yo.

\subsection*{Hipótesis críticas y limites}

\begin{enumerate}[leftmargin=*,nosep]
    \item \textbf{Atomicidad} de las lecturas/escrituras a \texttt{flag[]}
        y \texttt{turn}.
    \item \textbf{Modelo de memoria sequencially consistent.} Si la
        máquina puede reordenar las escrituras (1) y (2) ---lo que pasa
        en x86/ARM/RISC-V con buffers de escritura--- el algoritmo es
        semi-formalmente correcto pero prácticamente roto. Una
        implementación real de Peterson requiere \emph{barreras de
        memoria} (\texttt{MFENCE} en x86, \texttt{fence rw,rw} en
        RISC-V) entre los pasos (1) y (2).
\end{enumerate}

Por eso, en la práctica moderna, los kernels usan \emph{spinlocks}
construidos sobre instrucciones atómicas del hardware (TAS, CAS, FAA)
en lugar de Peterson. Peterson queda como demostración de existencia y
pieza pedagógica fundamental.

% =====================================================================
\section{Ejercicio 11 --- Productor-consumidor con semáforos POSIX}
% =====================================================================

Se pide implementar el problema del productor-consumidor con la API de
semáforos POSIX en dos variantes:

\begin{enumerate}[leftmargin=*,nosep]
    \item Usando POSIX threads.
    \item Usando procesos y \emph{named semaphores}, con buffer compartido
        (archivo o shared memory POSIX).
\end{enumerate}

\subsection*{Arquitectura común}

Solución clásica de Dijkstra con buffer acotado y tres semáforos:

\begin{center}
\small
\begin{tabular}{@{}llllp{4.5cm}@{}}
\toprule
Semáforo & Significado & Init & Lo decrementa & Lo incrementa \\
\midrule
\texttt{empty} & slots libres   & \texttt{BUF\_SIZE} & Productor antes de escribir & Consumidor tras leer \\
\texttt{full}  & slots ocupados & 0                  & Consumidor antes de leer    & Productor tras escribir \\
\texttt{mtx}   & mutex sobre el buffer & 1          & Quien va a tocar el buffer  & Quien lo libera \\
\bottomrule
\end{tabular}
\end{center}

\begin{lstlisting}[style=textstyle]
Productor                          Consumidor
---------                          ----------
sem_wait(empty)                    sem_wait(full)
sem_wait(mtx)                      sem_wait(mtx)
  buffer[in] = item                  item = buffer[out]
  in = (in+1) % BUF_SIZE             out = (out+1) % BUF_SIZE
sem_post(mtx)                      sem_post(mtx)
sem_post(full)                     sem_post(empty)
\end{lstlisting}

\paragraph{Justificación del orden de los \texttt{wait}.}

\begin{itemize}[leftmargin=*,nosep]
    \item \texttt{sem\_wait(empty)} antes que \texttt{sem\_wait(mtx)}.
        Si fuera al revés, el productor podría tomar el mutex teniendo
        el buffer lleno y dejar al consumidor sin chance de avanzar:
        \emph{deadlock}.
    \item El mutex se mantiene \textbf{sólo} para la actualización del
        buffer y de los índices. El \emph{trabajo de produzir/consumir}
        queda fuera para no serializar todo el pipeline.
    \item \texttt{empty} y \texttt{full} actúan como variables de
        condición: bloquean al productor cuando el buffer está lleno y
        al consumidor cuando está vacío.
\end{itemize}

\subsection*{11.1 Versión con POSIX threads}

Como los hilos comparten memoria, alcanza con declarar el buffer y los
semáforos como variables globales del proceso. Los semáforos son
\emph{sin nombre} (\texttt{sem\_t} con \texttt{sem\_init}).

\begin{lstlisting}[style=cstyle]
#include <pthread.h>
#include <semaphore.h>
#include <stdio.h>
#include <stdlib.h>
#include <unistd.h>

#define BUF_SIZE 8
#define N_ITEMS  32

static int buffer[BUF_SIZE];
static int in_idx  = 0;
static int out_idx = 0;

static sem_t empty, full, mtx;

static void *producer(void *arg) {
    (void)arg;
    for (int i = 0; i < N_ITEMS; i++) {
        int item = i * i;

        sem_wait(&empty);
        sem_wait(&mtx);
        buffer[in_idx] = item;
        in_idx = (in_idx + 1) % BUF_SIZE;
        printf("[P] produjo %d (in=%d)\n", item, in_idx);
        sem_post(&mtx);
        sem_post(&full);

        usleep(10 * 1000);
    }
    return NULL;
}

static void *consumer(void *arg) {
    (void)arg;
    for (int i = 0; i < N_ITEMS; i++) {
        sem_wait(&full);
        sem_wait(&mtx);
        int item = buffer[out_idx];
        out_idx = (out_idx + 1) % BUF_SIZE;
        printf("    [C] consumio %d (out=%d)\n", item, out_idx);
        sem_post(&mtx);
        sem_post(&empty);

        usleep(25 * 1000);
    }
    return NULL;
}

int main(void) {
    sem_init(&empty, 0, BUF_SIZE);
    sem_init(&full,  0, 0);
    sem_init(&mtx,   0, 1);

    pthread_t tp, tc;
    pthread_create(&tp, NULL, producer, NULL);
    pthread_create(&tc, NULL, consumer, NULL);
    pthread_join(tp, NULL);
    pthread_join(tc, NULL);

    sem_destroy(&empty); sem_destroy(&full); sem_destroy(&mtx);
    return 0;
}
\end{lstlisting}

\begin{lstlisting}
cc -O2 -Wall -Wextra -pthread -o prod_cons_threads prod_cons_threads.c
./prod_cons_threads
\end{lstlisting}

\begin{lstlisting}[style=textstyle]
[P] produjo 0 (in=1)
[P] produjo 1 (in=2)
    [C] consumio 0 (out=1)
[P] produjo 4 (in=3)
    [C] consumio 1 (out=2)
[P] produjo 9 (in=4)
...
\end{lstlisting}

El productor adelanta al consumidor hasta que se llena el buffer
(\texttt{empty} se vuelve 0); ahí se bloquea en
\texttt{sem\_wait(empty)} hasta que el consumidor libere un slot.

\paragraph{Nota.} En macOS/Darwin \texttt{sem\_init()} está deprecado y
devuelve \texttt{ENOSYS} en runtime; este programa está pensado para
GNU/Linux. Para macOS, ver la versión 11.2 con \emph{named semaphores}.

\subsection*{11.2 Versión con procesos y \emph{named semaphores}}

Los procesos no comparten memoria por defecto, así que hace falta:

\begin{enumerate}[leftmargin=*,nosep]
    \item Una región de \textbf{shared memory POSIX} (\texttt{shm\_open}
        + \texttt{ftruncate} + \texttt{mmap}) para el buffer y los
        índices.
    \item \textbf{Semáforos con nombre} (\texttt{sem\_open}) que vivan
        en el namespace global del kernel.
\end{enumerate}

\begin{lstlisting}[style=cstyle]
#include <fcntl.h>
#include <semaphore.h>
#include <stdio.h>
#include <stdlib.h>
#include <string.h>
#include <sys/mman.h>
#include <sys/stat.h>
#include <sys/wait.h>
#include <unistd.h>

#define BUF_SIZE 8
#define N_ITEMS  32

#define NAME_EMPTY "/so_pc_empty"
#define NAME_FULL  "/so_pc_full"
#define NAME_MTX   "/so_pc_mtx"
#define NAME_SHM   "/so_pc_shm"

typedef struct {
    int  buffer[BUF_SIZE];
    int  in_idx;
    int  out_idx;
} shared_t;

static void die(const char *m) { perror(m); exit(1); }

int main(void) {
    sem_unlink(NAME_EMPTY);
    sem_unlink(NAME_FULL);
    sem_unlink(NAME_MTX);
    shm_unlink(NAME_SHM);

    int fd = shm_open(NAME_SHM, O_CREAT | O_RDWR, 0600);
    if (fd < 0) die("shm_open");
    if (ftruncate(fd, sizeof(shared_t)) != 0) die("ftruncate");

    shared_t *sh = mmap(NULL, sizeof(shared_t),
                       PROT_READ | PROT_WRITE, MAP_SHARED, fd, 0);
    if (sh == MAP_FAILED) die("mmap");
    memset(sh, 0, sizeof(*sh));

    sem_t *empty = sem_open(NAME_EMPTY, O_CREAT | O_EXCL, 0600, BUF_SIZE);
    sem_t *full  = sem_open(NAME_FULL,  O_CREAT | O_EXCL, 0600, 0);
    sem_t *mtx   = sem_open(NAME_MTX,   O_CREAT | O_EXCL, 0600, 1);
    if (empty==SEM_FAILED || full==SEM_FAILED || mtx==SEM_FAILED)
        die("sem_open");

    pid_t pid = fork();
    if (pid < 0) die("fork");

    if (pid == 0) {
        for (int i = 0; i < N_ITEMS; i++) {
            sem_wait(full);
            sem_wait(mtx);
            int item = sh->buffer[sh->out_idx];
            sh->out_idx = (sh->out_idx + 1) % BUF_SIZE;
            printf("    [C pid=%d] consumio %d (out=%d)\n",
                   (int)getpid(), item, sh->out_idx);
            sem_post(mtx);
            sem_post(empty);
            usleep(25 * 1000);
        }
        exit(0);
    }

    for (int i = 0; i < N_ITEMS; i++) {
        int item = i * i;
        sem_wait(empty);
        sem_wait(mtx);
        sh->buffer[sh->in_idx] = item;
        sh->in_idx = (sh->in_idx + 1) % BUF_SIZE;
        printf("[P pid=%d] produjo %d (in=%d)\n",
               (int)getpid(), item, sh->in_idx);
        sem_post(mtx);
        sem_post(full);
        usleep(10 * 1000);
    }

    waitpid(pid, NULL, 0);

    sem_close(empty); sem_close(full); sem_close(mtx);
    sem_unlink(NAME_EMPTY);
    sem_unlink(NAME_FULL);
    sem_unlink(NAME_MTX);
    munmap(sh, sizeof(*sh));
    close(fd);
    shm_unlink(NAME_SHM);
    return 0;
}
\end{lstlisting}

Compilación:

\begin{lstlisting}
# Linux: hace falta -lrt para shm_open
cc -O2 -Wall -Wextra -o prod_cons_processes prod_cons_processes.c -lrt
# macOS: shm_open/sem_open viven en libc, no se necesita -lrt
cc -O2 -Wall -Wextra -o prod_cons_processes prod_cons_processes.c
\end{lstlisting}

\begin{lstlisting}[style=textstyle]
[P pid=12345] produjo 0 (in=1)
    [C pid=12346] consumio 0 (out=1)
[P pid=12345] produjo 1 (in=2)
[P pid=12345] produjo 4 (in=3)
    [C pid=12346] consumio 1 (out=2)
...
\end{lstlisting}

\texttt{pid} distinto en cada línea: confirma que productor y consumidor
son \emph{procesos separados} que sin embargo comparten el mismo
\emph{buffer} y los mismos semáforos.

\paragraph{Por qué \texttt{O\_EXCL} y \texttt{sem\_unlink}.}
\texttt{sem\_open} con \texttt{O\_CREAT | O\_EXCL} falla si el semáforo
ya existe; eso obliga a una creación limpia y avisa si una corrida
anterior dejó basura en el sistema.

\paragraph{Por qué shared memory y no archivo.} Es más rápido (no toca
el filesystem); \texttt{mmap MAP\_SHARED} sobre el \texttt{fd} de
\texttt{shm\_open} pone una sola estructura en RAM accesible a ambos
procesos sin sincronización adicional más allá del mutex POSIX. En
sistemas modernos las regiones aparecen en \texttt{/dev/shm/<nombre>} y
son fáciles de inspeccionar durante el desarrollo.

\subsection*{Por qué la solución no tiene los problemas clásicos}

\begin{itemize}[leftmargin=*,nosep]
    \item \textbf{Sin lost wakeup}: los semáforos contadores recuerdan
        los \texttt{post()} aunque no haya nadie esperando. Si el
        consumidor llega antes que el productor,
        \texttt{sem\_wait(full)} lo deja dormido; cuando el productor
        hace \texttt{sem\_post(full)}, el consumidor se despierta.
    \item \textbf{Sin deadlock}: ambos toman primero los contadores y
        después el mutex; el mutex se libera en cada iteración. No hay
        ciclos.
    \item \textbf{Sin starvation}: \texttt{sem\_wait} en POSIX es FIFO
        en la mayoría de las implementaciones (no es estrictamente
        requerido por el estándar pero es lo habitual en glibc/Linux).
    \item \textbf{Sin race condition} sobre \texttt{in\_idx}/
        \texttt{out\_idx}: ambos viven dentro de la región crítica.
\end{itemize}

\subsection*{Conexión con la teoría}

El algoritmo es exactamente el del Listado~8 de la teoría
(\emph{Concurrencia e IPC}, \emph{Semáforos}). El ejercicio cubre los
\emph{dos modelos} de concurrencia descritos en clase:

\begin{enumerate}[leftmargin=*,nosep]
    \item \textbf{Memoria compartida intra-proceso} (versión con
        threads): los hilos comparten todas las variables globales
        naturalmente.
    \item \textbf{Shared memory POSIX entre procesos} (versión con
        procesos): la región se crea explícitamente con
        \texttt{shm\_open} + \texttt{mmap}; los semáforos requieren ser
        \emph{named} para ser visibles entre namespaces de procesos
        distintos. Esto corresponde al mecanismo IPC \emph{Shared
        memory} de la Tabla~1 de las Notas (\emph{IPC en UNIX}).
\end{enumerate}

La elección de tres semáforos sigue el patrón clásico: dos contadores
para sincronización por condición (replazan al \texttt{wait}/\texttt{signal}
de las variables de condición) más un mutex binario para exclusión
mutua sobre el buffer.

% =====================================================================
\section{Conclusiones}
% =====================================================================

El práctico recorre, de manera práctica y progresiva, los tres pilares
del modelo de tareas del SO: la \emph{representación} del proceso
(estados, atributos, descriptores expuestos por \texttt{ps} y
\texttt{/proc}), su \emph{planificación} (interrupciones del timer,
context switches, prioridades con \texttt{nice}/\texttt{renice}) y su
\emph{coordinación} ante recursos compartidos (race conditions,
\texttt{flock}, mutexes POSIX, monitores Java, semáforos). Los dos
ejercicios teóricos sobre operaciones del kernel (8 y 9) muestran que
la diferencia esencial entre desalojo por \emph{quantum} y bloqueo por
I/O es el estado al que se transita: \texttt{RUNNABLE} en un caso,
\texttt{SLEEPING} en el otro. Esa misma distinción ---tiempo compartido
voluntario vs. espera por evento--- es la que justifica toda la
arquitectura de las \emph{wait queues}, los semáforos y las variables
de condición que se materializan en el productor-consumidor del
Ejercicio~11. Peterson, finalmente, hace el cierre teórico: la
exclusión mutua puede resolverse por software, pero las hipótesis de
atomicidad y consistencia secuencial que requiere muestran por qué los
sistemas reales recurren a instrucciones atómicas del hardware. Todo
el práctico converge en una sola idea: la concurrencia es \emph{no
determinismo controlado}, y el objetivo del SO es ofrecer abstracciones
suficientes para que ese no determinismo no altere los objetivos de
cómputo.

\end{document}
